\abstract{%
Children presenting to the SickKids pediatric emergency department wait twice: once for physician assessment, and again for diagnostic test results. The Machine Learning Medical Directive (MLMD) system eliminates the second wait for high-confidence cases by predicting required tests from triage-time features (Canadian Triage and Acuity Scale (CTAS) score, vital signs, chief complaint, and demographics) and authorizing them in parallel. Each missed prediction returns the patient to the slower, serialized workflow, making recall the primary performance target. Label noise constrains this recall. Tests ordered at external facilities go unrecorded in the SickKids electronic health record (EHR), causing genuine test-ordering events to appear as negatives and teaching the model spurious clinical absences.

We address this across 483,705 pediatric ED encounters spanning June 2018 to January 2026. At the data level, we replace the production regex imputation system with a clinically validated Knowledge Graph and GraphRAG retrieval pipeline that recovers 11,110 previously unrecorded test events across four diagnostic pathways. At the training level, we apply LiLAW (Lightweight Learnable Adaptive Weighting)~\cite{moturu_lilaw_2025}, a meta-learning framework that down-weights residual false negatives through bilevel optimization over three learned scalar parameters.

GraphRAG imputation recovered 11,110 previously unrecorded test events, achieving precision of 0.944 on the most semantically ambiguous pathway (Testicular Ultrasounds) against 0.048 for lexical matching. On tabular benchmarks with asymmetric noise matching the MLMD structure, LiLAW improved PR-AUC by 1.2--2.9 percentage points across 10--40\% noise rates, with gains concentrated at the high-precision operating point clinically relevant to triage deployment. The downstream ablation integrating both interventions with the full MLMD architecture is in progress; component-level validation establishes the expected performance envelope and confirms the noise correction mechanisms function as designed.%
}
